% ===== 原文 pp.232–244（PDF p.247–259）· 第 5 章 §5.5–5.6 · tr-06 =====
% [原文 p.232]
\begin{lemma}
设 $P_1,\ldots,P_k$ 一般。设 $\gamma$ 是 $M_\R\setminus\operatorname{Sing}(\mathfrak D)$ 中的路径，$\gamma(0)=Q$，$\gamma(1)=Q'$。则对 $\dim\rho=1$，$\rho\in\Sigma$，
\begin{equation}
  L^d_{i,\rho}(Q')-L^d_{i,\rho}(Q)
  =L^d_{i,\gamma,\{0\}\to\rho}-\sum_{\substack{\sigma\in\Sigma\\ \rho\subsetneq\sigma}}L^d_{i,\gamma,\rho\to\sigma} \tag{5.34}
\end{equation}
而对 $\dim\sigma=2$，$\sigma\in\Sigma$，
\begin{equation}
  L^d_{i,\sigma}(Q')-L^d_{i,\sigma}(Q)
  =\sum_{\substack{\rho\in\Sigma\\ \dim\rho=1\\ \rho\subsetneq\sigma}}L^d_{i,\gamma,\rho\to\sigma}. \tag{5.35}
\end{equation}
\end{lemma}

\begin{proof}
只需对 $\gamma$ 是如定义 5.52 那样、在 $M_\R\setminus\operatorname{Supp}(\mathfrak D)$ 的两个相邻连通分支 $C_1$ 与 $C_2$ 中连接 $Q$ 与 $Q'$ 的短路径证明即可。假设在时刻 $t_0$，$\gamma(t_0)\in\mathfrak d_1\cap\cdots\cap\mathfrak d_s$，其中 $\mathfrak d_1,\ldots,\mathfrak d_s\in\mathfrak D$ 是射线。当然 $\dim\mathfrak d_i\cap\mathfrak d_j=1$。于是对任意 $i$ 取 $n_{\mathfrak d}=n_{\mathfrak d_i}$，可以写
\[
  \begin{aligned}
    \theta_{\gamma,\mathfrak D}(z^m)
    &=z^m\prod_{i=1}^{s}f_{\mathfrak d_i}^{\langle n_{\mathfrak d},r(m)\rangle}\\
    &=z^m\prod_{i=1}^{s}\bigl(1+c_{\mathfrak d_i}\langle n_{\mathfrak d},r(m)\rangle z^{m_{\mathfrak d_i}}\bigr)\\
    &=z^m+\sum_{i=1}^{s}c_{\mathfrak d_i}\langle n_{\mathfrak d},r(m)\rangle z^{m+m_{\mathfrak d_i}}.
  \end{aligned}
\]
这里最后一个等号来自 $c_{\mathfrak d_i}c_{\mathfrak d_j}=0$（$i\ne j$）。而这是由 $P_1,\ldots,P_k$ 一般这一假设保证的。事实上，若 $c_{\mathfrak d_i}c_{\mathfrak d_j}\ne0$，则对应于 $\mathfrak d_i$ 与 $\mathfrak d_j$ 的 Maslov 指标零树 $h_i$ 与 $h_j$ 会满足 $I(h_i)\cap I(h_j)=\varnothing$。然而把指标在 $I(h_i)$ 中的标记点作一般扰动，可以在不扰动 $\mathfrak d_j$ 的情况下形变 $\mathfrak d_i$，从而 $\mathfrak d_i\cap\mathfrak d_j=\varnothing$。

% [Tutor Note: p.232 (5.34)(5.35) 是 $L^d_{i,\omega}$ 的跨墙公式：跨墙时 $L^d_{i,\omega}$ 的变化只由 $\rho\to\sigma$（维数升一）的穿越项给出。证明中关键的一般位置论证是 $c_{\mathfrak d_i}c_{\mathfrak d_j}=0$——若两棵 Maslov 指标零树满足 $I(h_i)\cap I(h_j)=\varnothing$，可独立扰动相应标记点使两条射线不再相交，故同一交点处两条射线的系数之积必为零，于是乘积展开只保留一阶项。]

现在由定理 5.35，
\[
  W_k(Q')=\theta_{\gamma,\mathfrak D}(W_k(Q)).
\]
利用展开式 (5.22) 以及 $W_0(Q)=y_0+\sum_{j=0}^{2}z^{t_j}$，有
\[
  \begin{aligned}
    &\exp\bigl((W_k(Q')-W_0(Q'))/\hbar\bigr)\\
    &\quad=\exp\bigl((\theta_{\gamma,\mathfrak D}(W_k(Q))-W_0(Q))/\hbar\bigr)\\
    &\quad=\theta_{\gamma,\mathfrak D}\bigl(\exp((W_k(Q)-W_0(Q))/\hbar)\bigr)
      \cdot\exp\bigl((\theta_{\gamma,\mathfrak D}(W_0(Q))-W_0(Q))/\hbar\bigr)\\
    &\quad=\theta_{\gamma,\mathfrak D}\Bigl(\sum_{(c,m,\nu)\in S_k(Q)}c\hbar^{-\nu}z^m\Bigr)
      \Bigl(1+\hbar^{-1}\sum_{\ell=1}^{s}\sum_{j=0}^{2}c_{\mathfrak d_\ell}\langle n_{\mathfrak d},m_j\rangle z^{m_{\mathfrak d_\ell}+t_j}\Bigr)\\
    &\quad=\exp\bigl((W_k(Q)-W_0(Q))/\hbar\bigr)\\
    &\qquad+\sum_{(c,\nu,m)\in S_k(Q)}\sum_{\ell=1}^{s}\Bigl(c_{\mathfrak d_\ell}c\hbar^{-\nu}\bigl(\langle n_{\mathfrak d},r(m)\rangle z^{m+m_{\mathfrak d_\ell}}\\
    &\qquad\qquad+\hbar^{-1}\sum_{j=0}^{2}\langle n_{\mathfrak d},m_j\rangle z^{m+m_{\mathfrak d_\ell}+t_j}\bigr)\Bigr).
  \end{aligned}
\]

% [原文 p.233]
我们如下解释这一点。对每个 $(c,\nu,m)\in S_k(Q)$ 与每个 $\ell$，考察四项
\[
  c_{\mathfrak d_\ell}c\hbar^{-\nu}\Bigl(\langle n_{\mathfrak d},r(m)\rangle z^{m+m_{\mathfrak d_\ell}}
  +\hbar^{-1}\sum_{j=0}^{2}\langle n_{\mathfrak d},m_j\rangle z^{m+m_{\mathfrak d_\ell}+t_j}\Bigr).
\]
这四项对 $L^d_i(Q')$ 贡献表达式
\[
  c_{\mathfrak d_\ell}c\hbar^{-(3d+\nu-|m+m_{\mathfrak d_\ell}|)}
  \Bigl(\langle n_{\mathfrak d},r(m)\rangle D_i(d,m+m_{\mathfrak d_\ell})
  +\sum_{j=0}^{2}\langle n_{\mathfrak d},m_j\rangle D_i(d,m+m_{\mathfrak d_\ell}+t_j)\Bigr).
\]
可以验证，这一总贡献事实上为零：或者直接但繁琐地从 $D_i$ 的公式验证，或者对 $f=c\hbar^{-(\nu-1)}z^m$、$\theta=\theta_{\gamma,\mathfrak D}$ 应用引理 5.40。

现在若 $m+m_{\mathfrak d_\ell}$ 与 $m+m_{\mathfrak d_\ell}+t_j$（$0\le j\le 2$）全都落在同一个 $\tilde\omega_{\mathfrak d}$ 中，则这些项对任何 $\tau\in\Sigma$（包括 $\tau=\omega$）都不产生对 $L^d_{i,\tau}(Q')$ 的总贡献。另一方面，若 $m+m_{\mathfrak d_\ell}$ 与 $m+m_{\mathfrak d_\ell}+t_j$（$j=0,1,2$）不全落在同一个 $\omega\in\Sigma$ 的 $\tilde\omega_{\mathfrak d}$ 中，则这些项可以对不同的 $L^d_{i,\omega}(Q')$ 作出贡献。这只可能在 $m+m_{\mathfrak d}\in\tilde\omega_{\mathfrak d}$ 而 $m+m_{\mathfrak d}+t_j\in\tilde\tau_{\mathfrak d}$（对某个 $j$，其中 $\omega\subsetneq\tau\in\Sigma$、$\dim\tau=\dim\omega+1$、$m_j\in\tau$、$m_j\notin\omega$）时发生。在此情形，$L^d_{i,\tau}(Q')-L^d_{i,\tau}(Q)$ 有一项形如 $cc_{\mathfrak d_\ell}\langle n_{\mathfrak d},m_j\rangle\hbar^{-(3d+\nu-|m+m_{\mathfrak d_\ell}|)}D_i(d,m+m_{\mathfrak d_\ell}+t_j)$ 的贡献。于是 $L^d_{i,\omega}(Q')-L^d_{i,\omega}(Q)$ 必须有来自同一项、但符号相反的贡献。这就给出了引理。
\end{proof}

现在可以利用 $L^d_{i,\omega}(Q)$ 的渐近行为以及上面的跨墙公式来改写所需的表达式：

\begin{lemma}
设 $\gamma_j$ 是连接 $Q$ 与 $Q+sm_j$（$s\gg0$）的直线路径。设 $\gamma_{j,j+1}$ 是以 $Q$ 为基点的环路：它先从 $Q$ 直线走到 $Q+sm_j$，再沿大圆弧走到 $Q+sm_{j+1}$，然后从 $Q+sm_{j+1}$ 直线回到 $Q$。这里下标取模 $3$，且 $\gamma_{j,j+1}$ 总是逆时针环路。令 $\sigma_{j,j+1}=\rho_j+\rho_{j+1}$，它是 $\Sigma$ 中的二维锥。则
\begin{equation}
  L^d_i(Q)=L^d_{i,\{0\}}(Q)-\sum_{j=0}^{2}L^d_{i,\gamma_j,\{0\}\to\rho_j}
  -\sum_{j=0}^{2}L^d_{i,\gamma_{j,j+1},\rho_{j+1}\to\sigma_{j,j+1}}. \tag{5.36}
\end{equation}
\end{lemma}

\begin{proof}
由引理 5.51，对任何满足 $\rho_j\subseteq\sigma$ 的 $\sigma\in\Sigma$ 都有 $L^d_{i,\sigma}(Q+sm_j)=0$。于是由 (5.34) 与 (5.35) 得
\[
  \begin{aligned}
    L^d_{i,\rho_j}(Q)&=-L^d_{i,\gamma_j,\{0\}\to\rho_j}
      +\sum_{\substack{\sigma\in\Sigma\\ \rho_j\subsetneq\sigma}}L^d_{i,\gamma_j,\rho_j\to\sigma},\\
    L^d_{i,\sigma_{j,j+1}}(Q)&=-\sum_{\substack{\rho\in\Sigma\\ \dim\rho=1\\ \rho\subsetneq\sigma_{j,j+1}}}L^d_{i,\gamma_j,\rho\to\sigma_{j,j+1}}.
  \end{aligned}
\]
注意在第二个等式中我们打破了对称性。

% [原文 p.234]
把来自各 $\rho_j$ 与各 $\sigma_{j_1,j_2}$ 的贡献加起来，由引理 5.50 我们看到
\[
  \begin{aligned}
    L^d_i(Q)-L^d_{i,\{0\}}(Q)
    &={}-\sum_{j=0}^{2}L^d_{i,\gamma_j,\{0\}\to\rho_j}\\
    &\quad-\bigl(L^d_{i,\gamma_0,\rho_1\to\sigma_{0,1}}-L^d_{i,\gamma_1,\rho_1\to\sigma_{0,1}}\bigr)\\
    &\quad-\bigl(L^d_{i,\gamma_1,\rho_2\to\sigma_{1,2}}-L^d_{i,\gamma_2,\rho_2\to\sigma_{1,2}}\bigr)\\
    &\quad-\bigl(L^d_{i,\gamma_2,\rho_0\to\sigma_{2,0}}-L^d_{i,\gamma_0,\rho_0\to\sigma_{2,0}}\bigr).
  \end{aligned}
\]
再由引理 5.51，来自大圆弧的 $L^d_{i,\gamma_{j,j+1},\rho_{j+1}\to\sigma_{j,j+1}}$ 的贡献为零。因此
\[
  L^d_i(Q)-L^d_{i,\{0\}}(Q)
  =-\sum_{j=0}^{2}L^d_{i,\gamma_j,\{0\}\to\rho_j}
  -\sum_{j=0}^{2}L^d_{i,\gamma_{j,j+1},\rho_{j+1}\to\sigma_{j,j+1}},
\]
这正是所需结果。
\end{proof}

\subsubsection{最后的步骤}

现在我们已经把 $L^d_i(Q)$ 整理成可用的形式。我们将发现，(5.36) 中的每一项在方程 (5.23)–(5.25) 中都有几何解释：$L^d_{i,\{0\}}(Q)$ 给出这些方程右端来自形如 (5.4) 意义下的 $\langle\cdots\rangle^{\mathrm{trop}}_{d,\{0\}}$ 的贡献，$-L^d_{i,\gamma_j,\{0\}\to\rho_j}$ 给出由 $\langle\cdots\rangle^{\mathrm{trop}}_{d,\rho_j}$ 而来的贡献，而 $-L^d_{i,\gamma_{j,j+1},\rho_{j+1}\to\sigma_{j,j+1}}$ 给出由 $\langle\cdots\rangle^{\mathrm{trop}}_{d,\sigma_{j,j+1}}$ 而来的贡献。

我们先处理最简单的一项 $L^d_{i,\{0\}}$。这一情形的想法非常简单：周期积分恰好告诉我们把以 $Q$ 为边界的热带圆盘粘接起来、得到一条以高价数顶点映到 $Q$ 的热带曲线的所有方式。

\begin{lemma}
有
\[
  \begin{aligned}
    L^d_{i,\{0\}}(Q)
    &=\delta_{0,d}\delta_{0,i}
      +\sum_{\nu\ge i}\sum_{\substack{I\subseteq\{1,\ldots,k\}\\ I=\{i_1,\ldots,i_{3d-2+i-\nu}\}\\ i_1<\cdots<i_{3d-2+i-\nu}}}
      \langle P_{i_1},\ldots,P_{i_{3d-2+i-\nu}},\psi^\nu S\rangle^{\mathrm{trop}}_{d,\{0\}}u_I\hbar^{-(\nu+2-i)},
  \end{aligned}
\]
其中 $S=Q$、以 $Q$ 为顶点的热带直线 $L$，或在 $i=0,1$ 与 $2$ 的情形分别取 $M_\R$。
\end{lemma}

\begin{proof}
若 $d=0$，则 $S_k$ 中唯一对 $L^d_{i,\{0\}}$ 有贡献的元素 $(c,\nu,m)$ 是 $(1,0,0)$，它对应于 $e^{(W_k(Q)-W_0(Q))/\hbar}$ 中的常数项。当 $i=0$ 时它贡献 $1$，否则贡献 $0$，于是得到项 $\delta_{0,d}\delta_{0,i}$。

现在设 $d\ne0$。取 $(c,\nu,m)\in S_k$，$m=\sum_{i=0}^{2}n_it_i$。则 $(c,\nu,m)$ 只在 $n_0,n_1,n_2\le d$ 时对 $L^d_{i,\{0\}}$ 有贡献。写
\[
  c\hbar^{-\nu}z^m=\hbar^{-\nu}\prod_{i=1}^{\nu}\operatorname{Mono}(h_i)
  =\hbar^{-\nu}\prod_{i=1}^{\nu}\operatorname{Mult}(h_i)z^{\Delta(h_i)}u_{I(h_i)},
\]
其中 $h_i:\Gamma'_i\to M_\R$ 是边界为 $Q$ 的 Maslov 指标二圆盘，$1\le i\le\nu$。令 $\Gamma$ 为把 $\Gamma'_1,\ldots,\Gamma'_\nu$ 的出射顶点 $V_{\mathrm{out},i}$ 等同为单个顶点 $V_{\mathrm{out}}$，然后以 $V_{\mathrm{out}}$ 为顶点添加 $(d-n_0)+(d-n_1)+(d-n_2)+1$ 条额外无界边所得的图。定义 $h:\Gamma\to M_\R$ 使其在每个子图 $\Gamma'_i\subseteq\Gamma$ 上等于 $h_i$。此外，对 $0\le i\le 2$，$h$ 把其中 $d-n_i$ 条新

% [原文 p.235]
无界边映到射线 $Q+\R_{\ge0}m_i$。最后一条无界边标记为 $x$，并被 $h$ 收缩。

我断言 $h$ 现在是平衡的热带曲线。事实上，除可能除 $V_{\mathrm{out}}$ 外，平衡条件在 $\Gamma$ 的所有顶点处显然成立。考虑以 $V_{\mathrm{out}}$ 为端点的各条边的切向量。若 $m^{\mathrm{prim}}(h_i)$ 是 $\Gamma'_i$ 的出射边 $h(E_{\mathrm{out},i})$ 上指向远离 $Q$ 方向的本原切向量，则把平衡条件对 $\Gamma'_i$ 中不含 $V_{\mathrm{out},i}$ 的所有顶点求和，得到
\[
  r(\Delta(h_i))=w(E_{\mathrm{out},i})m^{\mathrm{prim}}(h_i).
\]
此外还有 $d-n_i$ 条权为一的边以 $V_{\mathrm{out}}$ 为顶点、方向为 $m_i$。现在按假设，
\[
  \sum_{i=1}^{\nu}\Delta(h_i)+\sum_{j=0}^{2}(d-n_j)t_j=d(t_0+t_1+t_2),
\]
于是作用 $r$ 得到
\[
  \sum_{i=1}^{\nu}r(\Delta(h_i))+\sum_{j=0}^{2}(d-n_j)m_j=0.
\]
而这正是 $V_{\mathrm{out}}$ 处的平衡条件。

于是该项 $c\hbar^{-\nu}z^m$ 对 $L^d_{i,\{0\}}$ 的贡献是
\[
  \begin{aligned}
    &\hbar^{-(3d+\nu-n_0-n_1-n_2)}D_i(d,n_0,n_1,n_2)\prod_{i=1}^{\nu}\operatorname{Mult}(h_i)u_{I(h_i)}\\
    &\quad=\hbar^{-(3d+\nu-n_0-n_1-n_2)}u_{I(h)}\operatorname{Mult}^i_x(h)
      \prod_{\substack{V\in\Gamma^{[0]}\\ V\notin E_x}}\operatorname{Mult}_V(h),
  \end{aligned}
\]
这是比较 $\operatorname{Mult}^i_x(h)$ 与 $D_i(d,n_0,n_1,n_2)$ 的定义得到的。注意 $h$ 中顶点 $V_{\mathrm{out}}$ 的价数 $\operatorname{Val}(V_{\mathrm{out}})$ 为 $\nu+3d-(n_0+n_1+n_2)+1$。设 $I(h)=\{i_1,\ldots,i_{3d-2+i-\nu'}\}$ 对某个 $\nu'$。注意 $h$ 是粘接 $\operatorname{Val}(V_{\mathrm{out}})-1$ 个 Maslov 指标二圆盘得到的，我们看到
\[
  \begin{aligned}
    \operatorname{Val}(V_{\mathrm{out}})-1
    &=\sum_{i=1}^{\nu}\bigl(|\Delta(h_i)|-\#I(h_i)\bigr)+(d-n_1)+(d-n_2)+(d-n_3)\\
    &=3d-(3d-2+i-\nu')=\nu'+2-i.
  \end{aligned}
\]
于是曲线 $h$ 恰好贡献了定义 5.13 的 (1)(a)、(2)(b) 或 (3)(d) 所给出的、对
\begin{equation}
  \langle P_{i_1},\ldots,P_{i_{3d-2+i-\nu'}},\psi^{\nu'}S\rangle^{\mathrm{trop}}_{d,\{0\}}u_{I(h)}\hbar^{-(\nu'+2-i)} \tag{5.37}
\end{equation}
的正确贡献。反之，给定任何对 (5.37) 有贡献的曲线 $h$，由引理 5.12 (3) 可知 $h$ 是按上述方式把某一族以 $Q$ 为端点的 Maslov 指标二圆盘粘接起来得到的。现在这些 Maslov 指标二圆盘中的某些可能没有标记点：它们是方向为 $m_0,m_1,m_2$ 的射线。这类射线至多有 $d$ 条，故定义非负整数 $n_0,n_1,n_2$，使得 $d-n_i$ 是方向 $m_i$ 上出现的射线条数。其余的 Maslov 指标二圆盘 $h_1,\ldots,h_\nu$ 都带标记点，因此在 $\exp((W_k(Q)-W_0(Q))/\hbar)$ 中出现项 $\hbar^{-\nu}\prod_{i=1}^{\nu}\operatorname{Mono}(h_i)$。这一

% [原文 p.236]
项对 $L^d_{i,\{0\}}$ 的贡献与曲线 $h$ 对 (5.37) 的贡献相同，理由如上。
\end{proof}

% [Tutor Note: p.235 上式中第三个因子 $(d-n_3)$ 应为 $(d-n_0)$：$V_{\mathrm{out}}$ 的价数已由同页给出为 $\nu+3d-(n_0+n_1+n_2)+1$，且 $n_3$ 在全文中从未定义；只有用 $(d-n_0)+(d-n_1)+(d-n_2)$ 才能得到 $3d-(3d-2+i-\nu')$。此为原文笔误，照译以存真。]

接着是 $-L^d_{i,\gamma_j,\{0\}\to\rho_j}$ 的解释：

\begin{lemma}
有
\[
  -L^d_{i,\gamma_j,\{0\}\to\rho_j}
  =\sum_{\nu\ge i-1}\sum_{\substack{I\subseteq\{1,\ldots,k\}\\ I=\{i_1,\ldots,i_{3d-2+i-\nu}\}\\ i_1<\cdots<i_{3d-2+i-\nu}}}
  \langle P_{i_1},\ldots,P_{i_{3d-2+i-\nu}},\psi^\nu S\rangle^{\mathrm{trop}}_{d,\rho_j}u_I\hbar^{-(\nu+2-i)},
\]
其中 $S=Q$、$L$ 或 $M_\R$，分别对应 $i=0$、$i=1$ 与 $i=2$。这里照常 $L$ 是以 $Q$ 为顶点的热带直线。
\end{lemma}

\begin{proof}
对 $i=0$ 这是空的，因为两边都为零，故设 $i\ge1$。不失一般性，考虑 $L^d_{i,\gamma_0,\{0\}\to\rho_0}$。这个量是来自每个点 $P\in Q+(\rho_0\setminus\{0\})$ 的贡献之和，而该点是 $Q+\rho_0$ 与某条射线 $\mathfrak d\in D=D(\Sigma,P_1,\ldots,P_k)$ 的交点。写
\[
  f_{\mathfrak d}=1+c_{\mathfrak d}z^{m_{\mathfrak d}}.
\]
考虑只穿过 $\mathfrak d$ 的一小段 $\gamma\subseteq\gamma_0$ 对 $L^d_{i,\gamma_0,\{0\}\to\rho_0}$ 的贡献。设 $\gamma$ 从 $Q_1$ 走到 $Q_2$。现在 $\mathfrak d$ 对应于一条经过 $P$ 的 Maslov 指标零树，把它在 $P$ 处切开，得到一个以 $P$ 为边界的 Maslov 指标零圆盘 $h_1:\Gamma'_1\to M_\R$。于是
\[
  f_{\mathfrak d}=1+w_{\Gamma'_1}(E_{\mathrm{out},1})\operatorname{Mult}(h_1)z^{\Delta(h_1)}u_{I(h_1)}.
\]
此外，$\exp((W_k(Q_1)-W_0(Q_1))/\hbar)$ 中的项 $cz^m\hbar^{-\nu}$ 来自 $\nu$ 个不同的以 $Q_1$ 为边界的 Maslov 指标二圆盘，记作 $h_2,\ldots,h_{\nu+1}$（每个都至少有一个标记点），并且
\[
  cz^m\hbar^{-\nu}=\hbar^{-\nu}\prod_{i=2}^{\nu+1}\operatorname{Mult}(h_i)z^{\Delta(h_i)}u_{I(h_i)}.
\]
为使这一项对 $L^d_{i,\gamma,\{0\}\to\rho_0}$ 有贡献，$m+m_{\mathfrak d}=\sum_{i=1}^{\nu+1}\Delta(h_i)$ 必须形如 $dt_0+n_1t_1+n_2t_2$，其中 $n_1,n_2\le d$。设情况如此。圆盘 $h_2,\ldots,h_{\nu+1}$ 形变为以 $P$ 为边界的圆盘，我们仍把它们记作 $h_2,\ldots,h_{\nu+1}$。把这些圆盘写成 $h_i:\Gamma'_i\to M_\R$。每个 $\Gamma'_i$（$1\le i\le\nu+1$）都有一个顶点 $V_{\mathrm{out},i}$。

利用这些数据，我们可以如下构造一条真正的热带曲线。令 $\Gamma$ 为把 $\Gamma'_1,\ldots,\Gamma'_{\nu+1}$ 中的所有出射顶点 $V_{\mathrm{out},i}$ 等同、得到带有可区分顶点 $V_{\mathrm{out}}$ 的图，然后以 $V_{\mathrm{out}}$ 为顶点添加
\[
  (d-n_1)+(d-n_2)+1
\]
条额外无界边所得的图。然后定义 $h:\Gamma\to M_\R$ 使其在 $\Gamma'_i\subseteq\Gamma$ 上与 $h_i$ 一致。此外，$h$ 把前 $d-n_1$ 条新无界边映到 $P+\R_{\ge0}m_1$，把接下来 $d-n_2$ 条新无界边映到 $P+\R_{\ge0}m_2$，最后一条无界边被收缩并标记为 $x$。注意 $\Gamma$ 在 $V_{\mathrm{out}}$ 处的价数为 $\operatorname{Val}(V_{\mathrm{out}})=\nu+1+(d-n_1)+(d-n_2)+1$。于是我们得到参数化曲线 $h:\Gamma\to M_\R$，且 $h(x)=P$。平衡条件需在 $V_{\mathrm{out}}$ 处检验，但如同引理 5.55 的证明，事实
\[
  \sum_{i=1}^{\nu+1}\Delta(h_i)+(d-n_1)t_1+(d-n_2)t_2=d(t_0+t_1+t_2)
\]

% [原文 p.237]
表明平衡条件确实在 $V_{\mathrm{out}}$ 处成立。

该项对 $-L^d_{i,\gamma,\{0\}\to\rho_0}$ 的贡献是
\begin{equation}
  -\langle n_{\mathfrak d},m_0\rangle u_{I(h)}w_{\Gamma'_1}(E_{\mathrm{out},1})
  \Bigl(\prod_{i=1}^{\nu+1}\operatorname{Mult}(h_i)\Bigr)
  D_i(d,d+1,n_1,n_2)\hbar^{-(\nu+3d-(d+n_1+n_2))}. \tag{5.38}
\end{equation}
注意 $n_{\mathfrak d}$ 是本原的，零化 $r(m_{\mathfrak d})$，且必须在 $-m_0$ 上取正值。此外，所选同构 $\bigwedge^2M\cong\Z$ 把 $w(E_{1,\mathrm{out}})n_{\mathfrak d}$ 在相差一个符号的意义下等同于 $X_{r(m_{\mathfrak d})}$。于是按定义 5.13 取 $m(h_1)=r(m_{\mathfrak d})$，我们看到
\[
  -\langle n_{\mathfrak d},m_0\rangle w_{\Gamma'_1}(E_{\mathrm{out},1})=|m(h_1)\wedge m_0|.
\]
因此 (5.38) 等于
\[
  |m(h_1)\wedge m_0|u_{I(h)}D_i(d,d+1,n_1,n_2)
  \Bigl(\prod_{\substack{V\in\Gamma^{[0]}\\ V\notin E_x}}\operatorname{Mult}_V(h)\Bigr)
  \hbar^{-(\operatorname{Val}(V_{\mathrm{out}})-2)}.
\]
现在按定义 5.13，通过直接比较 $D_i$ 的定义有 $D_i(d,d+1,n_1,n_2)=\operatorname{Mult}^{i-1}_x(h)$。此外，若 $I(h)=\{i_1,\ldots,i_{3d-2+i-\nu'}\}$ 对某个 $\nu'$，我们看到，由于 $h$ 是把一个 Maslov 指标零圆盘粘接到 $\operatorname{Val}(V_{\mathrm{out}})-2$ 个 Maslov 指标二圆盘上得到的，
\[
  \begin{aligned}
    \operatorname{Val}(V_{\mathrm{out}})-2
    &=\sum_{i=1}^{\nu+1}\bigl(|\Delta(h_i)|-\#I(h_i)\bigr)+(d-n_1)+(d-n_2)\\
    &=3d-(3d-2+i-\nu')=\nu'+2-i.
  \end{aligned}
\]
因此按定义 5.13，所考虑的项对 $-L^d_{i,\gamma_0,\{0\}\to\rho_0}$ 的贡献与曲线 $h$ 对
\begin{equation}
  \langle P_{i_1},\ldots,P_{i_{3d-2+i-\nu'}},\psi^{\nu'}S\rangle^{\mathrm{trop}}_{d,\rho_0}u_{I(h)}\hbar^{-(\nu'+2-i)} \tag{5.39}
\end{equation}
的贡献恰好相同，这正是所需。

反之，给定任何对 $\langle P_{i_1},\ldots,P_{i_{3d-2+i-\nu'}},\psi^{\nu'}S\rangle^{\mathrm{trop}}_{d,\rho_0}$ 有贡献且满足 $h(E_x)=P\in Q+(\rho_0\setminus\{0\})$ 的曲线 $h$，引理 5.12 (2) 的步骤表明 $h$ 必定恰好按上述方式出现。事实上，该引理表明 $h$ 可以通过把若干 Maslov 指标二圆盘与一个以 $P$ 为边界的 Maslov 指标零圆盘粘接得到。由定义 5.13 的 (2)(a) 或 (3)(c) 的条件，无标记的 Maslov 指标二圆盘都不能平行于 $m_0$。因此可假设 $h$ 是把 Maslov 指标二圆盘 $h_2,\ldots,h_{\nu+1}$、一个 Maslov 指标零圆盘 $h_1$、方向 $m_1$ 上的 $d-n_1$ 条射线以及方向 $m_2$ 上的 $d-n_2$ 条射线粘接得到的。于是对 $Q+\rho_0$ 上 $P$ 附近的点 $Q_1$，$\hbar^{-\nu}\prod_{i=2}^{\nu+1}\operatorname{Mono}(h_i)$ 出现在 $\exp((W_k(Q_1)-W_0(Q_1))/\hbar)$ 中，因此我们看到该项对 $-L^d_{i,\gamma_j,\{0\}\to\rho_j}$ 的贡献与曲线 $h$ 对 (5.39) 的贡献相同，这正是所需。
\end{proof}

现在我们来到最困难的情形，即解释 $-L^d_{i,\gamma_{j,j+1},\rho_{j+1}\to\sigma_{j,j+1}}$。第一步是把计算局部化到 $D=D(\Sigma,P_1,\ldots,P_k)$ 的奇点上。

\begin{lemma}
对每个点 $P\in\operatorname{Sing}(\mathfrak D)$，令 $\gamma_P$ 是绕 $P$ 的小逆时针环路，小到不绕过 $\operatorname{Sing}(\mathfrak D)$ 的其它点。

% [原文 p.238]
则
\[
  L^d_{i,\gamma_{j,j+1},\rho_{j+1}\to\sigma_{j,j+1}}
  =\sum_{P\in\operatorname{Sing}(\mathfrak D)\cap(Q+\sigma_{j,j+1})}L^d_{i,\gamma_P,\rho_{j+1}\to\sigma_{j,j+1}}.
\]
\end{lemma}

\begin{proof}
若 $C_1$ 与 $C_2$ 是 $M_\R\setminus\operatorname{Supp}\mathfrak D$ 的两个连通分支的闭包且 $\dim C_1\cap C_2=1$，而 $\gamma$ 是从 $\operatorname{Int}(C_1)$ 进入 $\operatorname{Int}(C_2)$、只穿过 $\operatorname{Int}(C_1\cap C_2)$ 一次的短路径，则 $L^d_{i,\gamma,\rho_{j+1}\to\sigma_{j,j+1}}$ 与 $\gamma$ 及其端点无关。此外，把 $\gamma$ 的方向反过来会改变 $L^d_{i,\gamma,\rho_{j+1}\to\sigma_{j,j+1}}$ 的符号。令 $\cP$ 是 $\operatorname{Supp}(\mathfrak D)$ 在 $M_\R$ 上诱导的多面体分解，$\check\cP$ 是对偶胞腔分解。后者这样得到：取点 $v_\sigma$ 为 $\cP$ 的极大胞腔 $\sigma$ 的重心。则 $\check\cP$ 由零维胞腔 $\{v_\sigma\}$、端点 $v_\sigma$ 与 $v_{\sigma'}$ 的一维胞腔（当 $\dim\sigma\cap\sigma'=1$ 时），以及对 $P\in\operatorname{Sing}(\mathfrak D)$ 的、顶点为 $\{v_\sigma\mid P\in\sigma\}$ 的二维胞腔组成。我们可以取 $\gamma_P$ 为 $\check\cP$ 中对应于 $P$ 的二维胞腔的逆时针边界。于是显然：若对 $Q+\sigma_{j,j+1}$ 中所有奇点 $P$ 求和 $L^d_{i,\gamma_P,\rho_{j+1}\to\sigma_{j,j+1}}$，则遍历 $\check\cP$ 的端点 $v_\sigma,v_{\sigma'}$ 的边所得的贡献会出现两次、符号相反，只要 $\sigma\cap\sigma'$ 包含于 $\sigma_{j,j+1}$。唯一存活下来的贡献来自 $M_\R\setminus\operatorname{Sing}(\mathfrak D)$ 中同伦于 $\gamma_{j,j+1}$ 的某条路径的边，这就得到结论。
\end{proof}

\begin{lemma}
设 $P\in\operatorname{Sing}(\mathfrak D)\cap(Q+\sigma_{j,j+1})$，并假设 $P\notin\{P_1,\ldots,P_k\}$。则
\[
  -L^d_{i,\gamma_P,\rho_{j+1}\to\sigma_{j,j+1}}
  =\sum_{\nu\ge0}\sum_{h}\operatorname{Mult}(h)u_{I(h)}\hbar^{-(\nu+2-i)},
\]
其中和取遍对 $\langle P_{i_1},\ldots,P_{i_{3d-2+i-\nu}},\psi^\nu S\rangle^{\mathrm{trop}}_{d,\sigma_{j,j+1}}$（对满足 $h(E_x)=P$ 的各种 $i_1<\cdots<i_{3d-2+i-\nu}$）有贡献的所有曲线 $h$。
\end{lemma}

\begin{proof}
注意对 $i=0$ 或 $1$ 这是空的，因为两边都为零，故可设 $i=2$。为节省书写，令
\[
  L_{P,j}:=L^d_{2,\gamma_P,\rho_{j+1}\to\sigma_{j,j+1}}.
\]
固定 $P$ 附近的基点 $Q'$。考虑
\[
  \exp\bigl((W_k(Q')-W_0(Q'))/\hbar\bigr)
\]
中形如
\begin{equation}
  c\hbar^{-\nu}z^m=\hbar^{-\nu}\prod_{i=3}^{\nu+2}\bigl(\operatorname{Mult}(h_i)z^{\Delta(h_i)}u_{I(h_i)}\bigr) \tag{5.40}
\end{equation}
的项，其中 $h_i$ 是以 $Q'$ 为边界的 Maslov 指标二圆盘，但没有任何 $h_i$ 来自在 $P$ 附近拐弯的断线。因此该项出现在 $\exp((W_k(Q'')-W_0(Q''))/\hbar)$ 中，对 $P$ 的小开邻域内所有一般的 $Q''$ 都如此。

% [Tutor Note: p.238 引理 5.58 只对 $i=2$ 有内容（$i=0,1$ 时两边为零，原文称 “this is vacuous”）；(5.40) 的乘积从 $i=3$ 开始、共 $\nu+2$ 个圆盘，与引理 5.56 中的 $\prod_{i=2}^{\nu+1}$（共 $\nu$ 个圆盘）不同，原因是这里的 $h_1,h_2$ 分别是墙两侧对应的 Maslov 指标零圆盘、已在穿墙时被消耗。核对时注意各引理中局部编号的差异。]

假设这样的项 $c\hbar^{-\nu}z^m$ 在 $\gamma_P$ 穿过某条射线 $\mathfrak d\in\mathfrak D$（$P\in\mathfrak d$，$P\ne\operatorname{Init}(\mathfrak d)$）时对 $L_{P,j}$ 有贡献。但 $\gamma_P$ 两次、以相反方向穿过这条射线 $\mathfrak d$，故 $c\hbar^{-\nu}z^m$ 会对 $L_{P,j}$ 贡献两次，但符号相反。于是这些贡献抵消，对 $L_{P,j}$ 的总和没有贡献。

因此我们只需分析当 $\gamma_P$ 穿过满足 $\operatorname{Init}(\mathfrak d)=P$ 的射线 $\mathfrak d$ 时出现的贡献，或来自 (5.40) 型单项式（其中某些 $h_i$ 来自在 $P$ 附近确实拐弯的断线）的贡献。

% [原文 p.239]
回顾命题 5.28 的论证，我们可以假设把经过 $P$ 的 $D$ 的射线集分解为
\[
  \{\mathfrak d_1,\ldots,\mathfrak d_n\}\cup\bigcup_{j=1}^{m}\mathfrak D_i,
\]
其中 $\mathfrak d_1,\ldots,\mathfrak d_n$ 不以 $P$ 为初始点，且对应于树 $h$，使得对任何出射边经过 $P$ 的 Maslov 指标零树 $h'$ 都有 $I(h)\cap I(h')\ne\varnothing$。

% [Tutor Note: p.239 上式中并集的指标为 $j=1,\dots,m$ 而并的项写作 $\mathfrak D_i$；按上下文（下文“$\mathfrak D_i$ 有三个元素”“$i\ne j$ 时……”以及定义 5.52 的分解）应为 $\bigcup_{i=1}^{m}\mathfrak D_i$。此处为原文笔误，照译以存真。]另一方面，$\mathfrak D_i$ 有三个元素，由两条不以 $P$ 为初始点的射线与一条以 $P$ 为初始点的射线组成，后者对应于把前两条射线对应的两棵树粘接所得的树。最后，若 $h$ 对应于 $\mathfrak D_i$ 中的射线、$h'$ 对应于 $\mathfrak D_j$ 中的射线且 $i\ne j$，则 $I(h)\cap I(h')\ne\varnothing$。

于是，为分析对 $L_{i,\gamma_P,\rho_{j+1}\to\sigma_{j,j+1}}$ 的贡献，上述讨论表明我们可以假设恰好有三条射线 $\mathfrak d_1,\mathfrak d_2,\mathfrak d_3$ 经过 $P$，且 $\operatorname{Init}(\mathfrak d_1),\operatorname{Init}(\mathfrak d_2)\ne P$、$\operatorname{Init}(\mathfrak d_3)=P$。现在 $\mathfrak d_1,\mathfrak d_2$ 对应于经过 $P$ 的 Maslov 指标零树，把它们切开，得到边界为 $P$ 的 Maslov 指标零圆盘 $h_i:\Gamma'_i\to M_\R$（$i=1,2$），并且对 $i=1,2$，
\[
  f_{\mathfrak d_i}=1+w_{\Gamma'_i}(E_{\mathrm{out},i})\operatorname{Mult}(h_i)z^{\Delta(h_i)}u_{I(h_i)}.
\]
我们现在分析可能对 $L_{P,j}$ 有贡献的额外项 $c\hbar^{-\nu}z^m$ 如何出现。以下假设 $c\hbar^{-\nu}z^m$ 如 (5.40) 中那样，且对应于 $h_3,\ldots,h_{\nu+2}$ 的断线都不在 $\mathfrak d_1,\mathfrak d_2$ 或 $\mathfrak d_3$ 处拐弯。

写
\[
  m+\Delta(h_1)+\Delta(h_2)=\sum_{j=0}^{2}n_jt_j.
\]
我们有以下各种额外贡献的可能性：

(I) 当 $\gamma_P$ 穿过 $\mathfrak d_3$ 时，$c\hbar^{-\nu}z^m$ 可能对 $L_{P,j}$ 有贡献。此贡献只可能在 $n_{j+2}\le n_j=d<n_{j+1}$ 时出现。

(II) 穿过 $\mathfrak d_1$ 之后，形如（略去系数）$z^{m+\Delta(h_1)}$ 与 $z^{m+\Delta(h_1)+t_\ell}$（$\ell=0,1,2$）的新项可能出现在 $\exp((W_k-W_0)/\hbar)$ 中。于是当我们穿过 $\mathfrak d_2$ 时，这些新项可能对 $L_{P,j}$ 有贡献。注意 $z^{m+\Delta(h_1)}$ 只在 $n_{j+2}\le n_j=d<n_{j+1}$ 时于穿过 $\mathfrak d_2$ 时贡献。项 $z^{m+\Delta(h_1)+t_j}$ 只在 $n_{j+2}\le d$、$n_j=d-1$ 且 $d<n_{j+1}$ 时贡献。项 $z^{m+\Delta(h_1)+t_{j+1}}$ 只在 $n_{j+2}\le n_j=d\le n_{j+1}$ 时贡献。项 $z^{m+\Delta(h_1)+t_{j+2}}$ 只在 $n_{j+2}<n_j=d<n_{j+1}$ 时贡献。

(III) 穿过 $\mathfrak d_2$ 之后，形如（略去系数）$z^{m+\Delta(h_2)}$ 与 $z^{m+\Delta(h_2)+t_\ell}$（$\ell=0,1,2$）的新项可能出现在 $\exp((W_k-W_0)/\hbar)$ 中。于是当我们穿过 $\mathfrak d_1$ 时，这些新项可能对 $L_{P,j}$ 有贡献。注意 $z^{m+\Delta(h_2)}$ 只在 $n_{j+2}\le n_j=d<n_{j+1}$ 时于穿过 $\mathfrak d_1$ 时贡献。项 $z^{m+\Delta(h_2)+t_j}$ 只在 $n_{j+2}\le d$、$n_j=d-1$ 且 $d<n_{j+1}$ 时贡献。项 $z^{m+\Delta(h_2)+t_{j+1}}$ 只在 $n_{j+2}\le n_j=d\le n_{j+1}$ 时贡献。项 $z^{m+\Delta(h_2)+t_{j+2}}$ 只在 $n_{j+2}<n_j=d<n_{j+1}$ 时贡献。

当这些额外贡献出现时，有三种情形。

\textit{情形 (a).} $n_{j+2}\le n_j=d<n_{j+1}$。在此情形，（略去系数）$z^m$ 在 $\gamma_P$ 穿过 $\mathfrak d_3$ 时给出 (I) 型贡献，而 $z^{m+\Delta(h_i)}$、$z^{m+\Delta(h_i)+t_{j+1}}$ 或 $z^{m+\Delta(h_i)+t_{j+2}}$（若 $n_{j+2}<d$）可能在 $\gamma_P$ 穿过 $\mathfrak d_1$ 与 $\mathfrak d_2$ 时给出 (II) 型或 (III) 型贡献。注意当我们遍历环路 $\gamma_P$ 时，这些单项式对 $L^d_{i,\sigma_{j,j+1}}$ 的总改变

% [原文 p.240]
是它们对 $L^d_{i,\gamma_P,\rho_j\to\sigma_{j,j+1}}$ 与 $L^d_{i,\gamma_P,\rho_{j+1}\to\sigma_{j,j+1}}$ 的贡献之和。然而，由于 $\gamma_P$ 是环路，对 $L^d_{i,\sigma_{j,j+1}}$ 改变的总贡献必然为零；又因 $d<n_{j+1}$，这些单项式都不对 $L^d_{i,\gamma_P,\rho_j\to\sigma_{j,j+1}}$ 的任何改变有贡献。因此这些单项式对 $L^d_{i,\gamma_P,\rho_{j+1}\to\sigma_{j,j+1}}=L_{P,j}$ 的总贡献也为零。

\textit{情形 (b).} $n_{j+2}\le d$，$n_j=d-1$，$d<n_{j+1}$。在此情形只有项 $z^{m+\Delta(h_i)+t_j}$（$i=1,2$）可能贡献。然而与情形 (a) 同样的论证表明这些项的总贡献为零。

\textit{情形 (c).} $n_{j+2}\le n_j=n_{j+1}=d$。在此情形，对 $L_{P,j}$ 的贡献只来自形如 $z^{m+\Delta(h_i)+t_{j+1}}$ 的项。取 $n_{\mathfrak d_i}$ 使得在 $\gamma_P$ 第一次经过 $\mathfrak d_i$ 的时刻 $\tau_i$ 有 $\langle n_{\mathfrak d_i},\gamma'_P(\tau_i)\rangle<0$。通过交换 $\mathfrak d_1$ 与 $\mathfrak d_2$ 的标号并适当选取基点 $Q'$，我们可以假设 $\gamma_P$ 先经过 $\mathfrak d_1$ 再经过 $\mathfrak d_2$，且 $\langle n_{\mathfrak d_i},m_{j+1}\rangle\ge0$（$i=1,2$）。写 $f_{\mathfrak d_i}=1+c_{\mathfrak d_i}z^{m_{\mathfrak d_i}}$（$i=1,2$）。

则在 $\gamma_P$ 第一次穿过 $\mathfrak d_1$ 之后，项
\[
  \langle n_{\mathfrak d_1},m_{j+1}\rangle cc_{\mathfrak d_1}\hbar^{-(\nu+1)}z^{m+\Delta(h_1)+t_{j+1}}
\]
出现在 $\exp((W_k(\gamma_P(t))-W_0(\gamma_P(t)))/\hbar)$ 中（并在我们第二次穿过 $\mathfrak d_1$ 时消失），因此当 $\gamma_P$ 第一次穿过 $\mathfrak d_2$ 时，我们得到对 $L_{P,j}$ 的贡献
\[
  \begin{aligned}
    &\langle n_{\mathfrak d_2},m_j\rangle\langle n_{\mathfrak d_1},m_{j+1}\rangle cc_{\mathfrak d_1}c_{\mathfrak d_2}\cdot\\
    &\qquad\cdot D_2(d,m+\Delta(h_1)+\Delta(h_2)+t_j+t_{j+1})
      \hbar^{-(\nu+3d-|m+\Delta(h_1)+\Delta(h_2)|)}.
  \end{aligned}
\]
另一方面，项 $\langle n_{\mathfrak d_2},m_{j+1}\rangle cc_{\mathfrak d_2}\hbar^{-(\nu+1)}z^{m+\Delta(h_2)+t_{j+1}}$ 出现在 $\gamma_P$ 第一次穿过 $\mathfrak d_2$ 之后（并在我们第二次穿过 $\mathfrak d_2$ 时消失），因此当 $\gamma_P$ 第二次穿过 $\mathfrak d_1$ 时，我们得到对 $L_{P,j}$ 的贡献
\[
  \begin{aligned}
    &\langle-n_{\mathfrak d_1},m_j\rangle\langle n_{\mathfrak d_2},m_{j+1}\rangle cc_{\mathfrak d_1}c_{\mathfrak d_2}\cdot\\
    &\qquad\cdot D_2(d,m+\Delta(h_1)+\Delta(h_2)+t_j+t_{j+1})
      \hbar^{-(\nu+3d-|m+\Delta(h_1)+\Delta(h_2)|)}.
  \end{aligned}
\]
注意
\[
  \begin{aligned}
    &\langle n_{\mathfrak d_2},m_j\rangle\langle n_{\mathfrak d_1},m_{j+1}\rangle
    -\langle n_{\mathfrak d_1},m_j\rangle\langle n_{\mathfrak d_2},m_{j+1}\rangle\\
    &\qquad=-|n_{\mathfrak d_1}\wedge n_{\mathfrak d_2}|
    =-|m^{\mathrm{prim}}(h_1)\wedge m^{\mathrm{prim}}(h_2)|,
  \end{aligned}
\]
因为 $n_{\mathfrak d_1},n_{\mathfrak d_2}$ 构成 $N_\R$ 的正定向基，而 $m_j,m_{j+1}$ 构成 $M_\R$ 的正定向基。

现在 Maslov 指标二圆盘 $h_3,\ldots,h_{\nu+2}$ 形变为以 $P$ 为边界的圆盘，我们仍记作 $h_3,\ldots,h_{\nu+2}$。于是我们可以把圆盘 $h_1,\ldots,h_{\nu+2}$ 连同方向 $m_{j+2}$ 上 $d-n_{j+2}$ 份无标记点的 Maslov 指标二圆盘粘接起来。它们在各自的出射顶点处粘接，得到一个顶点 $V_{\mathrm{out}}$，我们再添加一条带标号 $x$ 的额外无界边 $E_x$，也接在顶点 $V_{\mathrm{out}}$ 上。这给出图 $\Gamma$，其在 $V_{\mathrm{out}}$ 处的价数为 $\operatorname{Val}(V_{\mathrm{out}})=\nu+3+d-n_{j+2}$。于是我们得到参数化曲线 $h:\Gamma\to M_\R$，且 $h(x)=P$。再次容易验证 $V_{\mathrm{out}}$ 处的平衡条件。

% [原文 p.241]
于是，从 $c\hbar^{-\nu}z^m$ 出发按上述方式分析所得的对 $-L_{P,j}$ 的总贡献为
\[
  \begin{aligned}
    &|m^{\mathrm{prim}}(h_1)\wedge m^{\mathrm{prim}}(h_2)|w_{\Gamma'_1}(E_{\mathrm{out},1})w_{\Gamma'_2}(E_{\mathrm{out},2})\cdot\\
    &\qquad\cdot\operatorname{Mult}^0_x(h)\Bigl(\prod_{\substack{V\in\Gamma^{[0]}\\ V\notin E_x}}\operatorname{Mult}_V(h)\Bigr)
      \hbar^{-(\nu+3d-(2d+n_{j+2}))}u_{I(h)}\\
    &\quad=|m(h_1)\wedge m(h_2)|\operatorname{Mult}^0_x(h)
      \Bigl(\prod_{\substack{V\in\Gamma^{[0]}\\ V\notin E_x}}\operatorname{Mult}_V(h)\Bigr)
      \hbar^{-(\operatorname{Val}(V_{\mathrm{out}})-3)}u_{I(h)}.
  \end{aligned}
\]
可以看出，若 $I(h)=\{i_1,\ldots,i_{3d-\nu'}\}$ 对某个 $\nu'$，则由于 $h$ 是把两个 Maslov 指标零圆盘与 $\operatorname{Val}(V_{\mathrm{out}})-3$ 个 Maslov 指标二圆盘粘接得到的，
\[
  \begin{aligned}
    \operatorname{Val}(V_{\mathrm{out}})-3
    &=\sum_{i=1}^{\nu+2}\bigl(|\Delta(h_i)|-\#I(h_i)\bigr)+d-n_{j+2}\\
    &=3d-(3d-\nu')=\nu'.
  \end{aligned}
\]
于是我们看到，上述分析所得的从 $c\hbar^{-\nu}z^m$ 到 $-L_{P,j}$ 的贡献的系数，恰好是 $h$ 对
\begin{equation}
  \langle P_{i_1},\ldots,P_{i_{3d-\nu'}},\psi^{\nu'}M_\R\rangle^{\mathrm{trop}}_{d,\sigma_{j,j+1}}u_{I(h)}\hbar^{-\nu'} \tag{5.41}
\end{equation}
的贡献，这正是所需。

反之，给定对 (5.41) 有贡献且 $h(x)=P$ 的 $h$，可以用引理 5.12 (1) 在 $P$ 处把它切开、分解为热带圆盘。于是我们看到 $h$ 恰好按上述方式出现，完全如同引理 5.55 与 5.56 的证明。因此 $-L_{P,j}$ 是来自满足 $h(E_x)=P$ 的映射对 (5.41) 的贡献。
\end{proof}

\begin{lemma}
设 $P\in\sigma_{j,j+1}\cap\operatorname{Sing}(\mathfrak D)$，并假设 $P=P_\ell$ 对某个 $\ell$。则
\[
  -L^d_{i,\gamma_P,\rho_{j+1}\to\sigma_{j,j+1}}
  =u_\ell\hbar\delta_{d,0}\delta_{2,i}
  +\sum_{\nu\ge0}\sum_{h}\operatorname{Mult}(h)u_{I(h)}\hbar^{-(\nu+2-i)},
\]
其中和取遍对 $\langle P_{i_1},\ldots,P_{i_{3d-2+i-\nu}},\psi^\nu S\rangle^{\mathrm{trop}}_{d,\sigma_{j,j+1}}$（对满足 $h(E_x)=P$ 的各种 $i_1<\cdots<i_{3d-2+i-\nu}$）有贡献的所有曲线 $h$。
\end{lemma}

\begin{proof}
同样可设 $i=2$，并写
\[
  L_{P,j}=L^d_{2,\gamma_P,\rho_{j+1}\to\sigma_{j,j+1}}.
\]
取 $P_\ell$ 附近的基点 $Q'$。由注 5.29，$\mathfrak D$ 中含 $P_\ell$ 的射线与以 $P_\ell$ 为边界且不以 $P_\ell$ 为标记点的 Maslov 指标二圆盘一一对应。当 $Q'$ 充分接近 $P_\ell$ 时，这些 Maslov 指标二圆盘形变为以 $Q'$ 为边界的圆盘，故以 $P_\ell$ 为边界且不以 $P_\ell$ 为标记点的 Maslov 指标二圆盘与以 $Q'$ 为边界且不以 $P_\ell$ 为标记点的 Maslov 指标二圆盘一一对应。

若我们关心 $\exp((W_k(Q')-W_0(Q'))/\hbar)$ 中可能对 $L_{P_\ell,j}$ 有贡献的项，只需看 $\exp((W_k(Q')-W_0(Q'))/\hbar)$ 中

% [原文 p.242]
不以 $u_\ell$ 为因子的那些项，因为任何带该因子的项在我们穿过经过 $P_\ell$ 的射线时都不会产生新项。于是考虑形如
\begin{equation}
  c\hbar^{-\nu}z^m=\hbar^{-\nu}\prod_{p=1}^{\nu}\operatorname{Mult}(h_p)z^{\Delta(h_p)}u_{I(h_p)} \tag{5.42}
\end{equation}
的项，其中每个以 $Q'$ 为边界的圆盘 $h_p$ 都不经过 $P_\ell$，因而对应于一个以 $P_\ell$ 为边界的圆盘，我们也把它写作 $h_p:\Gamma'_p\to M_\R$。

把这些圆盘延拓为树并标记 $P_\ell$，得到 Maslov 指标零树，对应 $\mathfrak D$ 中以 $P_\ell$ 为初始点的射线 $\mathfrak d_p$。此外，我们有 $\mathfrak D$ 中的射线 $\mathfrak c_p$（$p=0,1,2$），以 $P_\ell$ 为初始点，对应于以 $Q'$ 为边界、无标记点的三个 Maslov 指标二圆盘。它们不出现在 $W_k(Q')-W_0(Q')$ 中，故与各 $\mathfrak d_p$ 不同。

以下我们写 $m=\sum_{p=1}^{\nu}\Delta(h_p)=\sum_{j=0}^{2}n_jt_j$，并取 $n_{\mathfrak d_p}$ 与 $n_{\mathfrak c_p}$ 的符号使得当 $\gamma_P$ 穿过相应射线时它们在 $\gamma'_P$ 上取负值。注意由于 $\gamma_P$ 逆时针，若使用由标准定向给出的等同 $\bigwedge^2M\cong\Z$（即 $m_1\wedge m_2\mapsto1$），则 $w_{\Gamma'_p}(E_{\mathrm{out},p})n_{\mathfrak d_p}=X_{r(\Delta(h_p))}$。于是
\begin{equation}
  \sum_{p=1}^{\nu}w_{\Gamma'_p}(E_{\mathrm{out},p})n_{\mathfrak d_p}=X_{r(m)}. \tag{5.43}
\end{equation}
另一方面，$\langle n_{\mathfrak c_{j+1}},m_j\rangle=-1$，$\langle n_{\mathfrak c_{j+2}},m_j\rangle=1$。

% [Tutor Note: p.241 (5.43) 的定向约定：$\bigwedge^2M\cong\Z$ 取 $m_1\wedge m_2\mapsto1$，配合逆时针的 $\gamma_P$，得 $w_{\Gamma'_p}(E_{\mathrm{out},p})n_{\mathfrak d_p}=X_{r(\Delta(h_p))}$；而 $\langle n_{\mathfrak c_{j+1}},m_j\rangle=-1$、$\langle n_{\mathfrak c_{j+2}},m_j\rangle=1$ 由三条无标记射线 $\mathfrak c_0,\mathfrak c_1,\mathfrak c_2$ 的方向 $m_0,m_1,m_2$ 的定向直接给出。这两条在情形 (a)、(c) 的符号计算中反复使用。]

现在我们可以把该项 $c\hbar^{-\nu}z^m$ 看成按以下四种方式对 $L_{P,j}$ 产生贡献：

(I) $\gamma_P$ 穿过 $\mathfrak d_l$（$1\le l\le\nu$）。则 $\exp((W_k(Q')-W_0(Q'))/\hbar)$ 中的项
\[
  \hbar^{-(\nu-1)}\prod_{\substack{p=1\\ p\ne l}}^{\nu}\operatorname{Mult}(h_p)z^{\Delta(h_p)}u_{I(h_p)}
\]
在 $n_{j+2}\le n_j=d<n_{j+1}$ 时对 $L_{P,j}$ 有贡献，此时贡献为
\[
  \langle n_{\mathfrak d_l},m_j\rangle\Bigl(\prod_{p=1}^{\nu}\operatorname{Mult}(h_p)u_{I(h_p)}\Bigr)
  u_\ell w_{\Gamma'_l}(E_{\mathrm{out},l})D_2(d,m+t_j)\hbar^{-(\nu+3d-|m|-1)}.
\]
注意这样的贡献要求 $\nu>0$。

(II) $\gamma_P$ 穿过 $\mathfrak c_j$。若 $c\hbar^{-\nu}z^m$ 在 $\gamma_P$ 穿过 $\mathfrak c_j$ 时对 $L_{P,j}$ 有贡献，其贡献会涉及因子 $\langle n_{\mathfrak c_j},m_j\rangle=0$，故没有贡献。

(III) $\gamma_P$ 穿过 $\mathfrak c_{j+1}$。若 $n_{j+2}\le n_j=d\le n_{j+1}$，我们从 $c\hbar^{-\nu}z^m$ 得到贡献，此时贡献为
\[
  \langle n_{\mathfrak c_{j+1}},m_j\rangle\Bigl(\prod_{p=1}^{\nu}\operatorname{Mult}(h_p)u_{I(h_p)}\Bigr)
  u_\ell D_2(d,m+t_j+t_{j+1})\hbar^{-(\nu+3d-|m|-1)}.
\]

(IV) $\gamma_P$ 穿过 $\mathfrak c_{j+2}$。若 $n_{j+2}<n_j=d<n_{j+1}$，我们从 $c\hbar^{-\nu}z^m$ 得到贡献，此时得到
\[
  \langle n_{\mathfrak c_{j+2}},m_j\rangle\Bigl(\prod_{p=1}^{\nu}\operatorname{Mult}(h_p)u_{I(h_p)}\Bigr)
  u_\ell D_2(d,m+t_j+t_{j+2})\hbar^{-(\nu+3d-|m|-1)}.
\]

% [原文 p.243]
现在考虑三种情形。

\textit{情形 (a).} $n_{j+2}\le n_j=d<n_{j+1}$，$\nu>0$。在此情形，忽略公共因子
\[
  \hbar^{-(\nu+3d-|m|-1)}u_\ell\prod_{p=1}^{\nu}\operatorname{Mult}(h_p)u_{I(h_p)},
\]
利用 (5.43) 与引理 5.43 中 $D_2$ 的公式，总贡献为
\[
  \begin{aligned}
    &\Bigl\langle\sum_{l=1}^{\nu}w_{\Gamma'_l}(E_{\mathrm{out},l})n_{\mathfrak d_l},m_j\Bigr\rangle D_2(d,m+t_j)-D_2(d,m+t_j+t_{j+1})\\
    &\qquad+\begin{cases}
      D_2(d,m+t_j+t_{j+2}) & n_{j+2}<d,\\
      0 & n_{j+2}=d
    \end{cases}\\
    &\quad=r(m)\wedge m_j(-1)^{n_j+n_{j+1}+1}\frac{(n_j-d)!(n_{j+1}-d-1)!}{(d-n_{j+2})!}\\
    &\qquad-(-1)^{n_j+n_{j+1}+2}\frac{(n_j-d)!(n_{j+1}-d)!}{(d-n_{j+2})!}\\
    &\qquad+(-1)^{n_j+n_{j+1}+1}\frac{(n_j-d)!(n_{j+1}-d-1)!}{(d-n_{j+2})!}(d-n_{j+2})\\
    &\quad=\bigl((n_{j+2}-n_{j+1})+(n_{j+1}-d)+(d-n_{j+2})\bigr)\cdot\\
    &\qquad\qquad\cdot(-1)^{n_j+n_{j+1}+1}\frac{(n_j-d)!(n_{j+1}-d-1)!}{(d-n_{j+2})!}\\
    &\quad=0.
  \end{aligned}
\]
所以此情形对 $L_{P,j}$ 没有贡献。

% [Tutor Note: p.243 情形 (a) 的消去是整个 §5.5.5 中“阶乘相消”计算的模型：括号 $(n_{j+2}-n_{j+1})+(n_{j+1}-d)+(d-n_{j+2})=0$ 使总贡献为零。可用引理 5.43 中 $D_2$ 的显式阶乘表达式逐项核验三项的系数 $(-1)^{n_j+n_{j+1}+1},(-1)^{n_j+n_{j+1}+2},(-1)^{n_j+n_{j+1}+1}(d-n_{j+2})$ 与分母 $(d-n_{j+2})!$。]

\textit{情形 (b).} $n_{j+2}\le n_j=d=n_{j+1}$，$\nu>0$。在此情形我们只从 (III) 得到贡献。此时我们可以把圆盘 $h_1,\ldots,h_\nu$ 连同方向 $m_{j+2}$ 上 $d-n_{j+2}$ 份无标记点的 Maslov 指标二圆盘粘接起来。它们在各自的出射顶点处粘接，得到一个顶点 $V_{\mathrm{out}}$，我们对某个 $l$ 添加两条额外带标记的无界边 $E_x$ 与 $E_{p_l}$，都接在 $V_{\mathrm{out}}$ 上。这给出图 $\Gamma$，其在 $V_{\mathrm{out}}$ 处的价数为 $\operatorname{Val}(V_{\mathrm{out}})=\nu+(d-n_{j+2})+2=\nu+3d-|m|+2$。于是我们得到参数化曲线 $h:\Gamma\to M_\R$，且 $h(V_{\mathrm{out}})=h(x)=h(p_l)=P_\ell$。此情形中 (III) 对 $-L_{P,j}$ 的贡献经检验容易看出是
\begin{equation}
  \operatorname{Mult}^0_x(h)\Bigl(\prod_{\substack{V\in\Gamma^{[0]}\\ V\notin E_x}}\operatorname{Mult}_V(h)\Bigr)
  u_{I(h)}\hbar^{-(\operatorname{Val}(V_{\mathrm{out}})-3)}. \tag{5.44}
\end{equation}
设 $I(h)=\{i_1,\ldots,i_{3d-\nu'}\}$ 对某个 $\nu'$，注意由于我们添加了映到 $P_\ell$ 的带标记边 $E_{p_l}$，有 $\ell\in I(h)$。由于 $h$ 是把 $\operatorname{Val}(V_{\mathrm{out}})-2$ 个 Maslov 指标二圆盘粘接得到的，我们有
\[
  \begin{aligned}
    \operatorname{Val}(V_{\mathrm{out}})-2
    &=\sum_{i=1}^{\nu}\bigl(|\Delta(h_i)|-\#I(h_i)\bigr)+d-n_{j+2}\\
    &=3d-(3d-\nu'-1)=\nu'+1.
  \end{aligned}
\]

% [原文 p.244]
于是我们看到 (5.44) 恰好是 $h$ 对
\[
  \langle P_{i_1},\ldots,P_{i_{3d-\nu'}},\psi^{\nu'}M_\R\rangle^{\mathrm{trop}}_{d,\sigma_{j,j+1}}u_{I(h)}\hbar^{-\nu'}
\]
（来自定义 5.13 (3)(b)）的贡献。如同我们考虑过的其它情形，反之任何这样的曲线 $h$ 都可以通过在 $P$ 处切开曲线而给出正确的单项式 $c\hbar^{-\nu}z^m$。

\textit{情形 (c).} $\nu=0$。$S_k$ 中只有一个元素 $(c,\nu,m)$ 满足 $\nu=0$，即对应于常数单项式 $1$ 的 $(1,0,0)$。所以 $n_0=n_1=n_2=0$，且除非 $d=0$ 我们没有贡献。同样，对 $L_{P,j}$ 的这一贡献只来自 (III)，并且是
\[
  \langle n_{\mathfrak c_{j+1}},m_j\rangle u_\ell D_2(0,t_j+t_{j+1})\hbar=-u_\ell\hbar.
\]
这给出了 $-L_{P,j}$ 中其余所声称的项。
\end{proof}

\begin{lemma}
有
\[
  \begin{aligned}
    -L^d_{i,\gamma_{j,j+1},\rho_{j+1}\to\sigma_{j,j+1}}
    &=\sum_{\substack{\ell\ \text{s.t.}\\ P_\ell\in Q+\sigma_{j,j+1}}}u_\ell\hbar\delta_{0,d}\delta_{2,i}\\
    &\quad+\sum_{\nu\ge0}\sum_{\substack{I\subseteq\{1,\ldots,k\}\\ I=\{i_1,\ldots,i_{3d-2+i-\nu}\}\\ i_1<\cdots<i_{3d-2+i-\nu}}}
      \langle P_{i_1},\ldots,P_{i_{3d-2+i-\nu}},\psi^\nu S\rangle^{\mathrm{trop}}_{d,\sigma_{j,j+1}}u_I\hbar^{-(\nu+2-i)},
  \end{aligned}
\]
其中 $S=Q$、$L$ 或 $M_\R$，分别对应 $i=0$、$i=1$ 与 $i=2$。
\end{lemma}

\begin{proof}
这不过是把引理 5.57、5.58 与 5.59 的结果合在一起。
\end{proof}

最后我们有：

\begin{proof}[定理 5.15 与 5.18 的证明]
由引理 5.46，只需证明 (5.23)、(5.24) 与 (5.25)。然而这些现在可由引理 5.54、5.55、5.56 与 5.60 推出。
\end{proof}

\subsection{参考文献与延伸阅读}

本节的素材完全取自论文 [42]。最初关于用以定义 Landau--Ginzburg 势的 Maslov 指标二全纯圆盘的提议应归功于 Cho 与 Oh [16]；然而他们只以这种方式定义了非扰动势。Fukaya、Oh、Ohta 与 Ono 的工作 [26] 与 [23] 追循了与此处所述类似的想法，用的是全纯圆盘而非热带圆盘。用全纯圆盘的优点是它可以应用于所有环面簇，而热带几何很难看到环面簇边界上的全纯曲线。另一方面，代价是必须用到巨著 [24]、[25]，而且很难做任何计算。Chan 与 Leung 在 [14] 中引入了用热带几何研究 Landau--Ginzburg 势的想法；同样，他们只处理了非扰动的版本。最后，关于 $\PP^2$ 的引力后代不变量的纯热带、非镜像对称的方法，见 [77]。

% [Tutor Note: p.244 §5.6 的文献编号 [42]、[16]、[26]、[23]、[24]、[25]、[14]、[77] 均按原文参考文献表照录不译（文献表由 Lead 统一保留英文）；原文称 [24]、[25] 为 immense work（巨著）。本节“后代”沿用第 2 章 c02a 的译法（gravitational descendent invariants 译“引力后代不变量”）。]

% [Tutor Note: p.244 原文说热带几何“难以看到环面簇边界上的全纯曲线”，这是与全纯圆盘方法对比的要点：热带化只记录 $\mathrm{Log}$ 像，靠近边界除子（如 $\PP^2$ 的坐标线）的曲线在热带极限下被压到扇的边界上，需额外的对数/非阿基米德数据才能恢复；这正是 [42] 引入“墙穿越 + 断线”这套机器（本书 §5.4–5.5）的动机。]
